\chapter{Conclusions, Limitations and Future Work}
\thispagestyle{empty}

\section{Conclusions}
This project developed an EEG pipeline for the detection of two neurodevelopmental disorders,
IDD and ADHD, with one common processing chain, and evaluated it strictly on people the model had
never seen. The main conclusions are as follows.

\begin{enumerate}
\item \textbf{The proposed pipeline detects both disorders on unseen people.} SMVMD common
      modes, their channel covariances in the Riemannian tangent space and LightGBM reach
      \R{MainDStwoBal}\% balanced accuracy and
      \R{MainDStwoSubj} children correct on ADHD, and classify all 14 participants correctly
      in all three IDD settings.
\item \textbf{It improves substantially on the best-reported method and matches the best
      published subject-wise result.} The starting method reaches \R{StepADStwoBal}\% on unseen
      ADHD children. The proposed pipeline is comparable with the best published subject-wise
      result on the same data, a deep convolutional network, while needing no GPU, about
      \R{RtTotal}~ms per window on one CPU core and about \R{RtTrain}~s of training.
\item \textbf{The channel covariance carries the information.} Representations that describe
      how channels co-vary (connectivity, Riemannian tangent space) outperform per-channel
      features, and most of the pipeline's importance lies in channel pairs.
\item \textbf{Gradient-boosted trees suit the tangent-space features.} LightGBM gives the highest
      score for the proposed configuration and for the undecomposed covariance, and both
      boosting libraries are ahead of logistic regression on every representation, although the
      differences between classifiers are small.
\item \textbf{The two disorders use different modes.} The inner cross-validation chose all four
      modes, dominated by $\delta$ and $\theta$, for ADHD, and mainly the $\alpha$ mode for IDD.
\item \textbf{The adaptive decomposition does not add accuracy by itself.} Fixed band-pass
      filters and the undecomposed signal perform as well or one to two points better; SMVMD's
      role is to provide an interpretable, data-driven split of
      each window into the oscillations it contains.
\item \textbf{Evaluation on new people is essential.} The proposed pipeline itself would be
      reported at \R{GapDStwoLeaky}\% on ADHD under window-level cross-validation, against
      \R{GapDStwoSubj}\% on unseen children; accuracies of 95--100\% in the literature are
      not comparable with subject-wise results.
\end{enumerate}

\section{Limitations}
\begin{itemize}
\item \textbf{Medication confound.} The ADHD children were medicated and the controls were not,
      so the ADHD results cannot be attributed to ADHD alone.
\item \textbf{Small IDD sample.} DS-1 has 14 adult participants; one participant changes the
      subject accuracy by about 7 points, so even 14/14 is a weak estimate, and the groups may
      differ in recording conditions, artefacts and age as well as in the disorder.
\item \textbf{One dataset per disorder.} Each disorder is represented by a single dataset with its
      own device, task and population, so a single model distinguishing the two disorders could
      not be evaluated meaningfully.
\item \textbf{No external validation.} No independent cohort was available, and nothing here
      supports clinical use.
\item \textbf{Fixed band edges.} The common modes use the conventional EEG bands; edges adapted to
      the data were not explored.
\end{itemize}

\section{Future work}
\begin{itemize}
\item \textbf{Better data}: both groups recorded under the same protocol and device, matched in
      age, with medication status recorded and ideally a washout period; for IDD, a cohort of
      children.
\item \textbf{External validation} of the frozen pipeline on an independent ADHD cohort.
\item \textbf{Data-driven common modes}: band edges learned from the distribution of SMVMD
      centre frequencies in the training data, or a weighting of the modes instead of a hard
      assignment.
\item \textbf{Subject-invariant representations}: Riemannian re-centring or other alignment of
      each person's covariances, which address the differences between individuals that limit
      all methods.
\item \textbf{Faster decomposition}: SMVMD accounts for almost all of the computation; a
      filter bank derived from the training-set mode frequencies would make the pipeline
      suitable for real-time use on a headset.
\item \textbf{More disorders}: extending the pipeline to autism spectrum disorder under the same
      protocol.
\end{itemize}
